\section{Conclusions}
\label{sec:conclusion}

This paper describes a flexible method to infuse multiple intrinsic alignments models into simulated weak lensing  catalogues. These are based on a linear or quadratic coupling between the galaxy intrinsic shapes and the projected tidal field, computed directly from projected mass maps, thus are straight forward to implement in existing cosmic shear simulations where light-cone mass shells are saved to disk. The IA models are further refined by considering various tracer types, connecting mock galaxies and the underlying dark matter field,  either assuming no connection, a linear bias or an HOD-based non-linear bias. All six models (2 coupling $\times$ 3 tracer types) are studied and compared to the commonly used NLA and TATT models, highlighting regimes where the agreements are at times excellent, at times inconsistent. Discrepancies are expected as the TATT and our IA models are  based on slightly different physical assumptions, notably concerning the non-linear galaxy bias model, or the impact of neglecting the radial component of the tidal fields, which we report in detail. In particular, the low-redshift $II$ term shows the largest disagreement for all models involving non-linear biased tracers.

We quantify the impact of such differences with a full MCMC analysis, where we find that the cosmology is well recovered only for the NLA and TT models, whereas all other models, which involve non-linear bias terms that are not fully captured by the current implementation of the TATT model, tend to push the inferred value of $S_8$ towards lower values, and $\Omega_{\rm m}$ towards higher values. The IA parameters are often degenerate with one another, and well recovered only for the NLA and TT models. These serve as good examples for understanding cosmic shear data analyses in which the IA sector is not guaranteed to follow neither the NLA nor the TATT model, and the type of cosmological biases shown here could be present in existing results reported in the literature.


Finally, we study the impact of these different IA models on non-Gaussian weak lensing statistics, including the lensing PDF, peak count, integrated bispectrum and $M_{\rm ap}^3$. We find that both the coupling type and the bias model have large effects. We isolate the source-coupling effect and find that it produces significant modifications to most data vectors, especially at low redshifts. Overall, the different IA models leave distinct signatures that should be differentiable with the upcoming data, however a full demonstration of this will require  an MCMC analysis based on these non-Gaussian probes, which we leave for future work.

The results shown in this work were obtained at a fixed cosmology, and we expect the amplitude of the impact of both SC and IA to be cosmology dependent. Applying our methods on mock with varying cosmology is straightforward and will allow us to investigate the cosmology dependence of our finding.

The current work will be key for upcoming non-Gaussian cosmic shear analyses, where the choice of IA model used will impact the inferred cosmology. Statistical tools such as empirical model selection \citep{Campos_Samuroff_IA} or Bayesian model averaging \citet{Grandon2024} will likely be required to protect the inference from IA model mis-specifications. 

Future work will include also a split by colour/morphology to better reflect the fact that different galaxy populations are affected differently by IA. Ideally, direct IA measurements from spectroscopic surveys \citep{IA_direct} will allow us to set priors on some of out model parameters, which should reflect in tighter cosmological constraints, but this is complicated by the complex selection of source galaxies, hence a full calibration is yet to be found.